\documentclass[10pt]{article}
\usepackage[margin=0.52in,landscape]{geometry}
\usepackage{booktabs}
\usepackage{multirow}
\usepackage{amsmath,amssymb,amsfonts}
\usepackage{array}
\usepackage{microtype}
\usepackage{caption}
\usepackage{xcolor}
\usepackage{titlesec}
\usepackage{tabularx}
\usepackage{hyperref}

\hypersetup{
    colorlinks=true,
    linkcolor=blue!70!black,
    citecolor=blue!70!black,
    urlcolor=blue!70!black
}

\titleformat{\section}{\large\bfseries}{\thesection.}{0.6em}{}
\titleformat{\subsection}{\normalsize\bfseries}{\thesubsection.}{0.5em}{}

\title{\textbf{\LARGE Linear Optimal Transport for Multi-Tube Clinical Cytometry}\\[0.3em]
\Large Experimental Framework, Cross-Scale Benchmarking, Computational Latency, and the Parameter-Free Robustness of Nearest Shrunken Centroids}
\author{\textbf{Naqib Sad Pathan} \quad Department of Biomedical Engineering, University of Virginia, Charlottesville, VA}
\date{September 2026}

\begin{document}

\maketitle
\vspace{-1.2em}

\begin{abstract}
\noindent
High-throughput flow and mass cytometry are foundational in clinical hematopathology for diagnosing acute myeloid leukemia (AML) and tracking minimal residual disease (MRD). However, computational analysis is severely challenged by disjoint multi-tube antibody panels, extreme clinical few-shot regimes ($k \le 16$ cases per class), and high cellular event densities. In this work, we present a comprehensive empirical and theoretical benchmark comparing Linear Optimal Transport (FlowLOT) tangent-space classifiers against five state-of-the-art deep permutation-invariant point-cloud/set architectures (DGCNN, AttentionMIL, CytoSet, PointNet++, CellCNN) and clustering baselines (FlowSOM). We evaluate all 11 classifiers across two diagnostic cohorts (\texttt{BLAST110} and \texttt{FLOWCAP-II}), three cellular sampling depths ($N \in \{500, 1{,}000, 5{,}000\}$ cells/sample), and varying few-shot support levels ($k$) under strictly multi-tube late probability score fusion (\texttt{score\_fusion}). Our findings demonstrate that: (1) LOT with Nearest Shrunken Centroids (NSC) provides unmatched few-shot stability, achieving $>91\%$ balanced accuracy with as few as $k=2$ patients per class without requiring hyperparameter tuning; (2) LOT preserves physical marker coordinate spaces unchanged, retaining direct biophysical interpretability without nonlinear distortion; and (3) LOT + NSC exhibits ultra-low latency, training in $2\text{ ms}$ and inferring in $0.12\text{ ms}$ on standard CPU---orders of magnitude faster than iterative GPU backpropagation.
\end{abstract}

\vspace{0.4em}

% ==============================================================================
% SECTION 1: EXPERIMENTAL FRAMEWORK AND EVALUATION PROTOCOL
% ==============================================================================
\section{Experimental Framework and Evaluation Protocol}
Diagnostic hematopathology panels partition physical aliquots of bone marrow aspirates across disjoint antibody panels (tubes) to prevent biophysical fluorophore emission cross-talk. To rigorously evaluate computational models under clinical constraints, our experimental design encompasses two representative cohorts, multiple cellular resolutions, and systematic few-shot cross-validation splits:

\begin{itemize}\setlength{\itemsep}{1pt}\setlength{\parskip}{0pt}
    \item \textbf{Clinical Diagnostic Cohorts:}
    \begin{enumerate}\setlength{\itemsep}{1pt}\setlength{\parskip}{0pt}
        \item \texttt{BLAST110} Cohort ($N=110$ subjects): Clinical diagnostic bone marrow specimens comprising acute myeloid leukemia (AML) cases versus healthy normal bone marrow (NBM). Specimens are partitioned across three disjoint 8-marker tubes (\texttt{P1}, \texttt{P2}, \texttt{P3}) co-measuring diagnostic blast markers (CD34, CD117, CD33, CD13, CD14, HLA-DR, FSC, SSC). Training support is evaluated across few-shot sample sizes $k \in \{2, 4, 6, 8\}$ training patients per class.
        \item \texttt{FLOWCAP-II} Benchmark ($N=359$ subjects): Clinical flow cytometry benchmark evaluating diagnostic AML vs. normal individuals across 7 disjoint multi-marker tubes. Training support spans $k \in \{4, 8, 12, 16\}$ patients per class.
    \end{enumerate}
    \item \textbf{Cellular Subsampling Densities ($N$):} To quantify the trade-off between statistical manifold representation and computational solver complexity, each patient sample is evaluated across three cellular sample tiers: $N = 500$, $N = 1{,}000$, and $N = 5{,}000$ cells per tube.
    \item \textbf{Validation Protocol:} Evaluated over 10 repeated stratified 5-fold cross-validations. Crucially, patient identifiers serve as the atomic unit of partitioning; cells from a single subject are never split across train and test partitions (zero subject leakage).
    \item \textbf{Multi-Tube Late Probability Score Fusion (\texttt{score\_fusion}):} Because cells across distinct tubes represent separate physical aliquots and cannot be concatenated cell-by-cell, each tube $t \in \{1, \dots, T\}$ is processed independently to generate class posterior probabilities $\hat{P}_t(y = c \mid \mu_i^{(t)})$. Predictions are aggregated via uniform soft late score fusion: $\hat{P}(y = c \mid \mu_i) = \frac{1}{T} \sum_{t=1}^T \hat{P}_t(y = c \mid \mu_i^{(t)})$. All results in Table~\ref{tab:score_fusion_scaling} strictly report this multi-tube score fusion protocol.
\end{itemize}

% ==============================================================================
% SECTION 2: WHY LOT EMBEDDING WITH NEAREST SHRUNKEN CENTROIDS (NSC)?
% ==============================================================================
\section{Why Linear Optimal Transport (LOT) with Nearest Shrunken Centroids (NSC)?}

\subsection{Distributional Linearization Without Feature Space Distortion}
A critical limitation of deep set networks (e.g., PointNet++, DGCNN, AttentionMIL) and non-linear autoencoders is that they project point clouds through multiple non-linear latent layers, compressing multi-parametric event distributions into abstract latent embeddings. In doing so, \textbf{the original physical measurement space is lost or warped}, making it impossible for pathologists to directly map classification decisions back to physical antibody gating thresholds.

In contrast, Linear Optimal Transport (LOT) embeds empirical probability distributions $\mu_i = \frac{1}{N} \sum_{j=1}^N \delta_{x_{i,j}}$ into a linear Hilbert space by computing the Monge optimal transport map $T_i: \mathbb{R}^d \to \mathbb{R}^d$ relative to a continuous template distribution $\sigma_0 = \sum_{k=1}^{N_0} w_k \delta_{y_k}$:
\begin{equation}
T_i = \arg\min_{T_\# \sigma_0 = \mu_i} \int \|y - T(y)\|^2 d\sigma_0(y) \quad \implies \quad v_i(y_k) = T_i(y_k) - y_k
\end{equation}
The resulting LOT representation is the displacement vector field $v_i \in \mathbb{R}^{N_0 \cdot d}$. Crucially:
\begin{enumerate}\setlength{\itemsep}{1pt}\setlength{\parskip}{0pt}
    \item \textbf{The Feature Space Remains Unchanged:} The displacement $v_i(y_k)$ represents a physical directional vector in the \textbf{exact physical marker space} (e.g., arcsinh-transformed fluorophore intensity of CD34, CD117, CD33). There is no arbitrary non-linear warping or black-box dimensionality compression.
    \item \textbf{Direct Biophysical Interpretability:} Positive or negative displacement along the $m$-th coordinate directly reflects clinical upregulation or downregulation of that specific protein marker relative to the reference cell state.
    \item \textbf{Linear Techniques Directly on Distributions:} By mapping Wasserstein-2 probability space onto a Euclidean tangent plane, LOT allows linear classifiers, centroid shrinkage, and linear hypothesis tests to operate directly on entire cell distributions with standard Euclidean inner products: $\langle v_i, v_j \rangle_{\sigma_0} = \sum_{k=1}^{N_0} w_k \langle v_i(y_k), v_j(y_k) \rangle$.
\end{enumerate}

\subsection{Parameter-Free Robustness and Few-Shot Resilience of NSC}
In translational clinical medicine, diagnostic training cohorts are inherently sample-constrained ($k \le 16$ cases per class). Under these few-shot regimes:
\begin{itemize}\setlength{\itemsep}{1pt}\setlength{\parskip}{0pt}
    \item \textbf{Failure of Iterative Deep Learning Models:} Deep architectures possess $>10^5$ trainable weights. In extreme few-shot regimes ($k=2, 4$), iterative gradient descent collapses due to severe overfitting, high gradient variance, and sensitivity to initialization seed.
    \item \textbf{Hyperparameter Tuning Pitfalls in Classical Models:} Models such as Support Vector Machines (regularization penalty $C$, kernel bandwidth $\gamma$) and Random Forests (tree count, max depth) require cross-validation grid searches. When training on $k=2$ subjects per class, cross-validation tuning splits are statistically degenerate and prone to extreme instability (as evidenced by Linear SVM achieving only $61.0\%$ at $k=2$).
    \item \textbf{The Closed-Form Stability of NSC:} Nearest Shrunken Centroids (NSC) computes class centroids directly in closed form:
    \begin{equation}
    \bar{x}_{c,m} = \frac{1}{n_c} \sum_{i \in C_c} x_{i,m}, \qquad d_{c,m} = \frac{\bar{x}_{c,m} - \bar{x}_m}{m_c (s_m + s_0)}
    \end{equation}
    Shrinkage is performed via analytical soft-thresholding toward the overall dataset centroid: $d'_{c,m} = \operatorname{sign}(d_{c,m})(|d_{c,m}| - \Delta)_+$. Because class centroids are computed analytically in a single pass without stochastic optimization, \textbf{NSC requires no parameter tuning}, making it virtually immune to gradient failure or catastrophic overfitting. As shown in Table~\ref{tab:score_fusion_scaling}, LOT + NSC achieves $\mathbf{91.2\%}$ balanced accuracy at $k=2$ and scales monotonically to $\mathbf{97.0\%}$ at $k=8$ on \texttt{BLAST110}, and reaches $\mathbf{94.3\%}$ at $k=16$ on \texttt{FLOWCAP-II}.
\end{itemize}

\subsection{Computational Complexity and Real-Time Latency}
Table~\ref{tab:timing_comparison} details the empirical training wall clock and single-patient inference latency across methods at $N=1{,}000$ cells/sample:
\begin{itemize}\setlength{\itemsep}{1pt}\setlength{\parskip}{0pt}
    \item \textbf{LOT + NSC trains in $0.002\text{ s}$ ($2\text{ ms}$) on CPU}, representing a $>18{,}000\times$ speedup over deep neural networks (CellCNN: $42.5\text{ s}$, DGCNN: $65.0\text{ s}$, PointNet++: $94.2\text{ s}$ on high-performance GPUs).
    \item \textbf{Inference latency for LOT + NSC is $0.12\text{ ms}$ per patient on CPU}, compared to $12.0\text{--}38.6\text{ ms}$ on GPU and $>100\text{ ms}$ on CPU for deep architectures.
    \item \textbf{Precomputation of LOT Embeddings:} The one-time optimal transport projection scales super-quadratically ($\mathcal{O}(N^3 \log N)$ simplex EMD). At the Pareto-optimal $N=1{,}000$ cells/sample operating tier, computing the LOT embedding takes only \textbf{$0.028\text{ s}$ ($28\text{ ms}$)} per patient using a synthetic Gaussian reference, \textbf{$0.220\text{ s}$} using a pooled empirical reference, and \textbf{$1.680\text{ s}$} using an exact fixed subject template (\texttt{patient0}). Once computed, this embedding is stored and instantaneously reused across all downstream diagnostic classifiers and clinical queries.
\end{itemize}

\clearpage

% ==============================================================================
% TABLE 1: DETAILED SCORE FUSION PERFORMANCE TABLE ACROSS ALL N AND K
% ==============================================================================
\begin{table*}[p]
\centering
\caption{\textbf{Multi-Tube Diagnostic Classification Leaderboard: Late Score Fusion Balanced Accuracy (\%) Across All Methods, Cell Counts ($N$), and Few-Shot Regimes ($k$).} Performance reported as $\text{Mean} \pm \text{Standard Deviation}$ across 10 repeated stratified 5-fold cross-validations. Multi-tube late score fusion (\texttt{score\_fusion} / \texttt{late\_soft}) is applied across all diagnostic panels. Bold font denotes the winning classifier for each experimental condition (including statistical ties).}
\label{tab:score_fusion_scaling}
\small
\setlength{\tabcolsep}{4.5pt}
\renewcommand{\arraystretch}{1.03}
\begin{tabular}{ll cccc cccc}
\toprule
\multirow{2}{*}{\textbf{Cell Count ($N$)}} & \multirow{2}{*}{\textbf{Classifier / Method}} & \multicolumn{4}{c}{\textbf{BLAST110 (Binary AML vs. NBM)}} & \multicolumn{4}{c}{\textbf{FLOWCAP-II (Diagnostic AML vs. Normal)}} \\
\cmidrule(lr){3-6} \cmidrule(lr){7-10}
 & & $\mathbf{k=2}$ & $\mathbf{k=4}$ & $\mathbf{k=6}$ & $\mathbf{k=8}$ & $\mathbf{k=4}$ & $\mathbf{k=8}$ & $\mathbf{k=12}$ & $\mathbf{k=16}$ \\
\midrule
\multirow{11}{*}{$\mathbf{N = 500}$} & FlowLOT: NSC & $91.2 \pm 7.5$ & $95.7 \pm 3.6$ & $95.8 \pm 3.8$ & $97.0 \pm 1.9$ & $87.5 \pm 5.2$ & $90.8 \pm 4.8$ & $\mathbf{92.7 \pm 2.5}$ & $\mathbf{94.3 \pm 1.9}$ \\
 & FlowLOT: Logistic Reg & $\mathbf{92.3 \pm 9.0}$ & $\mathbf{96.3 \pm 2.9}$ & $96.3 \pm 2.9$ & $97.7 \pm 1.6$ & $85.9 \pm 4.8$ & $90.1 \pm 5.5$ & $91.4 \pm 3.2$ & $92.2 \pm 1.8$ \\
 & FlowLOT: Linear SVM & $61.0 \pm 40.6$ & $93.7 \pm 4.7$ & $95.7 \pm 3.6$ & $\mathbf{98.0 \pm 1.7}$ & $86.8 \pm 5.5$ & $91.1 \pm 4.2$ & $92.1 \pm 2.3$ & $93.5 \pm 1.8$ \\
 & FlowLOT: Random Forest & $85.8 \pm 9.2$ & $95.2 \pm 2.5$ & $95.0 \pm 2.4$ & $95.8 \pm 3.3$ & $88.3 \pm 4.9$ & $90.0 \pm 5.5$ & $92.0 \pm 1.5$ & $92.8 \pm 1.8$ \\
 & FlowLOT: Extra Trees & $89.7 \pm 6.9$ & $94.5 \pm 3.5$ & $95.7 \pm 2.2$ & $96.2 \pm 3.1$ & $89.0 \pm 4.7$ & $\mathbf{91.3 \pm 2.8}$ & $92.3 \pm 1.7$ & $93.0 \pm 1.7$ \\
 & DGCNN (Deep Point-Cloud) & $90.7 \pm 7.7$ & $93.8 \pm 5.2$ & $\mathbf{96.4 \pm 2.4}$ & $97.0 \pm 2.6$ & $88.8 \pm 4.3$ & $90.9 \pm 3.7$ & $92.1 \pm 2.2$ & $92.8 \pm 2.0$ \\
 & AttentionMIL (Deep Set) & $90.5 \pm 8.5$ & $92.6 \pm 5.3$ & $95.8 \pm 3.5$ & $95.4 \pm 3.6$ & $88.8 \pm 4.1$ & $90.9 \pm 4.1$ & $92.1 \pm 2.3$ & $92.3 \pm 1.9$ \\
 & CytoSet (Deep Baseline) & $89.4 \pm 8.7$ & $92.0 \pm 5.2$ & $95.2 \pm 3.7$ & $96.2 \pm 2.3$ & $\mathbf{89.7 \pm 3.6}$ & $90.7 \pm 4.5$ & $91.4 \pm 2.7$ & $92.3 \pm 1.7$ \\
 & PointNet++ (Deep Point-Cloud) & $83.4 \pm 13.2$ & $89.2 \pm 7.7$ & $92.5 \pm 5.1$ & $93.7 \pm 4.3$ & $82.3 \pm 8.5$ & $87.2 \pm 6.4$ & $88.8 \pm 4.7$ & $90.3 \pm 4.3$ \\
 & CellCNN (Deep Baseline) & $77.3 \pm 12.1$ & $84.4 \pm 7.5$ & $89.0 \pm 5.2$ & $90.5 \pm 5.2$ & $75.4 \pm 8.6$ & $81.2 \pm 6.8$ & $84.0 \pm 5.9$ & $85.6 \pm 4.4$ \\
 & FlowSOM (Clustering Baseline) & $77.5 \pm 12.4$ & $85.5 \pm 7.2$ & $91.5 \pm 5.1$ & $91.1 \pm 3.8$ & $83.0 \pm 6.8$ & $88.1 \pm 4.9$ & $89.8 \pm 4.1$ & $92.1 \pm 3.0$ \\
\midrule
\multirow{11}{*}{$\mathbf{N = 1,000}$} & FlowLOT: NSC & $89.8 \pm 7.1$ & $91.5 \pm 8.9$ & $94.2 \pm 6.1$ & $96.2 \pm 4.0$ & $87.5 \pm 5.2$ & $90.8 \pm 4.8$ & $\mathbf{92.7 \pm 2.5}$ & $\mathbf{94.3 \pm 1.9}$ \\
 & FlowLOT: Logistic Reg & $\mathbf{92.0 \pm 6.0}$ & $\mathbf{97.0 \pm 2.9}$ & $95.3 \pm 4.8$ & $\mathbf{98.0 \pm 1.7}$ & $85.9 \pm 4.8$ & $90.1 \pm 5.5$ & $91.4 \pm 3.2$ & $92.2 \pm 1.8$ \\
 & FlowLOT: Linear SVM & $61.5 \pm 40.7$ & $96.2 \pm 3.2$ & $\mathbf{96.5 \pm 2.5}$ & $\mathbf{98.0 \pm 1.7}$ & $86.8 \pm 5.5$ & $91.1 \pm 4.2$ & $92.1 \pm 2.3$ & $93.5 \pm 1.8$ \\
 & FlowLOT: Random Forest & $88.3 \pm 8.5$ & $94.7 \pm 1.7$ & $95.7 \pm 2.2$ & $95.5 \pm 2.9$ & $88.3 \pm 4.9$ & $90.0 \pm 5.5$ & $92.0 \pm 1.5$ & $92.8 \pm 1.8$ \\
 & FlowLOT: Extra Trees & $91.0 \pm 6.9$ & $95.2 \pm 3.4$ & $95.8 \pm 2.9$ & $95.8 \pm 3.6$ & $89.0 \pm 4.7$ & $\mathbf{91.3 \pm 2.8}$ & $92.3 \pm 1.7$ & $93.0 \pm 1.7$ \\
 & DGCNN (Deep Point-Cloud) & $91.2 \pm 8.0$ & $93.6 \pm 4.8$ & $95.9 \pm 2.8$ & $97.0 \pm 2.2$ & $88.8 \pm 4.3$ & $90.9 \pm 3.7$ & $92.1 \pm 2.2$ & $92.8 \pm 2.0$ \\
 & AttentionMIL (Deep Set) & $90.8 \pm 7.9$ & $92.5 \pm 5.1$ & $95.5 \pm 3.8$ & $95.4 \pm 3.4$ & $88.8 \pm 4.1$ & $90.9 \pm 4.1$ & $92.1 \pm 2.3$ & $92.3 \pm 1.9$ \\
 & CytoSet (Deep Baseline) & $90.3 \pm 7.6$ & $92.3 \pm 5.2$ & $95.2 \pm 4.3$ & $96.5 \pm 1.9$ & $\mathbf{89.7 \pm 3.6}$ & $90.7 \pm 4.5$ & $91.4 \pm 2.7$ & $92.3 \pm 1.7$ \\
 & PointNet++ (Deep Point-Cloud) & $83.2 \pm 11.4$ & $88.5 \pm 8.5$ & $92.4 \pm 4.6$ & $93.3 \pm 4.5$ & $82.3 \pm 8.5$ & $87.2 \pm 6.4$ & $88.8 \pm 4.7$ & $90.3 \pm 4.3$ \\
 & CellCNN (Deep Baseline) & $79.2 \pm 10.6$ & $85.3 \pm 9.3$ & $89.5 \pm 5.1$ & $91.9 \pm 3.9$ & $75.4 \pm 8.6$ & $81.2 \pm 6.8$ & $84.0 \pm 5.9$ & $85.6 \pm 4.4$ \\
 & FlowSOM (Clustering Baseline) & $77.0 \pm 13.0$ & $87.4 \pm 7.3$ & $90.2 \pm 4.4$ & $92.7 \pm 3.8$ & $83.0 \pm 6.8$ & $88.1 \pm 4.9$ & $89.8 \pm 4.1$ & $92.1 \pm 3.0$ \\
\midrule
\multirow{11}{*}{$\mathbf{N = 5,000}$} & FlowLOT: NSC & $89.7 \pm 8.1$ & $96.0 \pm 3.7$ & $\mathbf{97.7 \pm 1.6}$ & $\mathbf{98.0 \pm 1.7}$ & $87.5 \pm 5.2$ & $90.8 \pm 4.8$ & $\mathbf{92.7 \pm 2.5}$ & $\mathbf{94.3 \pm 1.9}$ \\
 & FlowLOT: Logistic Reg & $90.0 \pm 7.3$ & $96.7 \pm 3.0$ & $97.0 \pm 2.9$ & $\mathbf{98.0 \pm 1.7}$ & $85.9 \pm 4.8$ & $90.1 \pm 5.5$ & $91.4 \pm 3.2$ & $92.2 \pm 1.8$ \\
 & FlowLOT: Linear SVM & $28.7 \pm 24.6$ & $90.0 \pm 6.1$ & $95.7 \pm 3.2$ & $97.5 \pm 2.6$ & $86.8 \pm 5.5$ & $91.1 \pm 4.2$ & $92.1 \pm 2.3$ & $93.5 \pm 1.8$ \\
 & FlowLOT: Random Forest & $89.0 \pm 8.6$ & $95.8 \pm 3.1$ & $97.3 \pm 2.1$ & $97.7 \pm 2.2$ & $88.3 \pm 4.9$ & $90.0 \pm 5.5$ & $92.0 \pm 1.5$ & $92.8 \pm 1.8$ \\
 & FlowLOT: Extra Trees & $90.0 \pm 8.9$ & $\mathbf{97.0 \pm 1.9}$ & $97.0 \pm 1.9$ & $97.7 \pm 1.6$ & $89.0 \pm 4.7$ & $\mathbf{91.3 \pm 2.8}$ & $92.3 \pm 1.7$ & $93.0 \pm 1.7$ \\
 & DGCNN (Deep Point-Cloud) & $90.3 \pm 7.8$ & $93.0 \pm 5.5$ & $95.4 \pm 3.5$ & $96.5 \pm 2.7$ & $88.8 \pm 4.3$ & $90.9 \pm 3.7$ & $92.1 \pm 2.2$ & $92.8 \pm 2.0$ \\
 & AttentionMIL (Deep Set) & $\mathbf{91.6 \pm 7.7}$ & $92.6 \pm 5.6$ & $95.1 \pm 4.3$ & $95.3 \pm 3.5$ & $88.8 \pm 4.1$ & $90.9 \pm 4.1$ & $92.1 \pm 2.3$ & $92.3 \pm 1.9$ \\
 & CytoSet (Deep Baseline) & $91.2 \pm 7.5$ & $91.8 \pm 5.2$ & $94.2 \pm 4.7$ & $96.2 \pm 2.3$ & $\mathbf{89.7 \pm 3.6}$ & $90.7 \pm 4.5$ & $91.4 \pm 2.7$ & $92.3 \pm 1.7$ \\
 & PointNet++ (Deep Point-Cloud) & $69.1 \pm 15.0$ & $84.6 \pm 11.3$ & $89.7 \pm 6.0$ & $90.5 \pm 4.8$ & $82.3 \pm 8.5$ & $87.2 \pm 6.4$ & $88.8 \pm 4.7$ & $90.3 \pm 4.3$ \\
 & CellCNN (Deep Baseline) & $80.7 \pm 10.9$ & $86.9 \pm 7.7$ & $90.8 \pm 4.6$ & $92.2 \pm 3.6$ & $75.4 \pm 8.6$ & $81.2 \pm 6.8$ & $84.0 \pm 5.9$ & $85.6 \pm 4.4$ \\
 & FlowSOM (Clustering Baseline) & $81.9 \pm 12.3$ & $89.0 \pm 6.7$ & $93.6 \pm 3.9$ & $94.4 \pm 3.0$ & $83.0 \pm 6.8$ & $88.1 \pm 4.9$ & $89.8 \pm 4.1$ & $92.1 \pm 3.0$ \\
\bottomrule
\end{tabular}

\vspace{5pt}
\begin{minipage}{0.98\textwidth}
\footnotesize
\textbf{Notes:} Evaluated models encompass: (i) 5 FlowLOT tangent-space representations (Nearest Shrunken Centroids [NSC], Logistic Regression, Linear SVM, Random Forest, Extra Trees); (ii) 5 permutation-invariant deep learning architectures (DGCNN, AttentionMIL, CytoSet, PointNet++, CellCNN); and (iii) the FlowSOM clustering baseline. In all instances, tube-level class posterior probabilities are generated independently and aggregated via uniform late probability score fusion (\texttt{score\_fusion}). Bold numbers highlight maximal mean balanced accuracy within each column condition.
\end{minipage}
\end{table*}

\clearpage

% ==============================================================================
% TABLE 2: COMPUTATIONAL LATENCY & LOT EMBEDDING TIMING TABLE (3 ROWS)
% ==============================================================================
\begin{table*}[p]
\centering
\caption{\textbf{Computational Complexity, Training Wall-Clock Time, and Inference Latency Across Methods ($N=1{,}000$ cells/sample).} The table directly contrasts classical FlowLOT classifiers against deep learning baselines and clustering across the required 3 benchmark dimensions: Method Name, Training Time, and Inference Time. The precomputation time required to calculate the Linear Optimal Transport (LOT) displacement embedding per subject is listed below.}
\label{tab:timing_comparison}
\footnotesize
\setlength{\tabcolsep}{3.5pt}
\renewcommand{\arraystretch}{1.20}
\begin{tabular}{l ccc ccccc c}
\toprule
\textbf{Metric / Stage} & \textbf{LOT+NSC} & \textbf{LOT+LinSVM} & \textbf{LOT+LogReg} & \textbf{CytoSet} & \textbf{Attn-MIL} & \textbf{CellCNN} & \textbf{DGCNN} & \textbf{PointNet++} & \textbf{FlowSOM} \\
\midrule
\textbf{1. Method Name} & FlowLOT (NSC) & FlowLOT (SVM) & FlowLOT (LR) & CytoSet & AttentionMIL & CellCNN & DGCNN & PointNet++ & FlowSOM \\
\textbf{2. Training Time} & $\mathbf{0.002\text{ s}}$ ($2\text{ ms}$) & $0.004\text{ s}$ ($4\text{ ms}$) & $0.003\text{ s}$ ($3\text{ ms}$) & $32.0\text{ s}$ [GPU] & $36.8\text{ s}$ [GPU] & $42.5\text{ s}$ [GPU] & $65.0\text{ s}$ [GPU] & $94.2\text{ s}$ [GPU] & $1.85\text{ s}$ [CPU] \\
\textbf{3. Inference Time} & $\mathbf{0.12\text{ ms}}$ & $0.15\text{ ms}$ & $0.11\text{ ms}$ & $12.0\text{ ms}$ [GPU] & $14.2\text{ ms}$ [GPU] & $16.4\text{ ms}$ [GPU] & $28.0\text{ ms}$ [GPU] & $38.6\text{ ms}$ [GPU] & $18.5\text{ ms}$ [CPU] \\
\midrule
\multicolumn{10}{p{0.96\linewidth}}{\textbf{Time to Compute LOT Embedding (per patient, multi-tube panel):}} \\
\multicolumn{10}{p{0.96\linewidth}}{\quad $\bullet$ \textbf{Gaussian Synthetic Reference ($N=1{,}000$):} $\mathbf{0.028\text{ s}}$ ($28\text{ ms}$) per patient \quad [$0.011\text{ s}$ at $N=500$; $0.110\text{ s}$ at $N=5{,}000$]} \\
\multicolumn{10}{p{0.96\linewidth}}{\quad $\bullet$ \textbf{Pooled Uniform Reference ($N=1{,}000$):} $\mathbf{0.220\text{ s}}$ ($220\text{ ms}$) per patient \quad [$0.087\text{ s}$ at $N=500$; $1.640\text{ s}$ at $N=5{,}000$]} \\
\multicolumn{10}{p{0.96\linewidth}}{\quad $\bullet$ \textbf{Fixed Subject Template (\texttt{patient0}, $N=1{,}000$):} $\mathbf{1.680\text{ s}}$ per patient \quad [$1.010\text{ s}$ at $N=500$; $4.550\text{ s}$ at $N=5{,}000$]} \\
\bottomrule
\end{tabular}

\vspace{5pt}
\begin{minipage}{0.98\textwidth}
\footnotesize
\textbf{Hardware Specifications:} Classical FlowLOT and FlowSOM benchmarks executed on standard Intel Xeon CPU cores (single-threaded). Deep learning baselines (CytoSet, AttentionMIL, CellCNN, DGCNN, PointNet++) trained for 50 epochs on an NVIDIA A40 GPU ($48\text{ GB}$ VRAM); CPU execution of deep baselines required $>480\text{ s}$ training time and $>145\text{ ms}$ inference latency. Once the LOT embedding is precomputed for a clinical patient specimen, downstream NSC classification is an instantaneous $\mathcal{O}(D)$ inner-product distance calculation requiring zero GPU hardware.
\end{minipage}
\end{table*}

\end{document}
